\section{Introduction}

The hard problem of consciousness (Chalmers, 1995) highlights the explanatory gap between physical processes and subjective experience. While various approaches have been proposed, from denying consciousness exists at all to treating it as a separate property alongside physical properties, the relationship between experience and physical reality remains contested. This paper explores a panpsychist approach, building on recent developments in the field (Strawson, 2006; Goff, 2017) and drawing inspiration from process philosophy (Whitehead, 1929) and neutral monism (Russell, 1927). It also integrates insights from the author's direct experience in meditation, systems architecture, data modelling, and software development, bringing practical computational perspectives to these philosophical questions.

This paper presents a framework that takes experience as ontologically primitive, similar to how Integrated Information Theory (Tononi, 2008) treats consciousness as fundamental, but with a different mathematical structure. This framework, which is termed ``Vibe Theory,'' should be understood as a philosophical exploration rather than a scientific hypothesis ready for empirical testing.

The overarching goal of this framework is to develop a unified vocabulary and terminology that can describe all phenomena in the universe, from quantum mechanics to consciousness to spiritual experience, within a single coherent language. While currently philosophical in nature, the ultimate aim is to refine these concepts toward mathematical rigor, enabling formal modeling and the generation of testable hypotheses. This work represents an initial step toward a more normalized and standardized scientific language that bridges the gap between subjective experience and objective reality.

This work is a philosophical and conceptual exploration rather than a scientific theory. It does not claim to derive, reproduce, or replace established physical laws, nor does it offer empirical predictions at this stage. References to concepts from physics, computation, and mathematics are intended as abstract analogies or sources of structural inspiration, not as assertions of physical equivalence or explanatory reduction. Any apparent parallels to existing scientific theories should be read as interpretive or metaphorical unless explicitly stated otherwise. The framework is presented as a candidate ontology and vocabulary for interdisciplinary discussion, open to revision, critique, and formal refinement.
